\chapter{2002年大纲卷（文）}
\section{单选题}
\begin{problem}
直线 \((1 + a)x + y + 1 = 0\) 与圆 \(x^{2} + y^{2} - 2x = 0\) 相切，则 \(a\) 的值为（\quad）

\choices
    {\(\pm 1\)}
    {\(\pm 2\)}
    {\(1\)}
    {\(-1\)}
\end{problem}

\begin{answer}
D
\end{answer}

\begin{solution}
圆 \(x^{2}+y^{2}-2x=0\) 可化为 \((x-1)^{2}+y^{2}=1\)，所以圆心为 \(C(1,0)\)，半径为 \(1\). 直线与圆相切的充要条件是圆心到直线的距离等于半径，因此
\[
\frac{|(1+a)\cdot1+1|}{\sqrt{(1+a)^{2}+1}}=1.
\]
两边平方并化简，得 \((a+2)^{2}=(a+1)^{2}+1\)，即 \(2a+2=0\)，所以 \(a=-1\)，故选 D.
\end{solution}

\begin{problem}
复数 \(\left(\frac{1}{2} + \frac{\sqrt{3}}{2}\i\right)^3\) 的值是（\quad）

\choices
    {\(-\i\)}
    {\(\i\)}
    {\(-1\)}
    {\(1\)}
\end{problem}

\begin{answer}
C
\end{answer}

\begin{solution}
注意到 \(\frac{1}{2}+\frac{\sqrt{3}}{2}\i=\cos\frac{\pi}{3}+\i\sin\frac{\pi}{3}\). 由棣莫弗定理，
\[
\left(\cos\frac{\pi}{3}+\i\sin\frac{\pi}{3}\right)^{3}=\cos\pi+\i\sin\pi=-1.
\]
所以原复数的值为 \(-1\)，故选 C.
\end{solution}

\begin{problem}
不等式 \((1 + x)(1 - |x|) > 0\) 的解集是（\quad）

\choices
    {\(\{x\mid 0\leqslant x < 1\}\)}
    {\(\{x\mid x < 0 \text{ 且 } x\neq -1\}\)}
    {\(\{x\mid -1 < x < 1\}\)}
    {\(\{x\mid x < 1 \text{ 且 } x\neq -1\}\)}
\end{problem}

\begin{answer}
D
\end{answer}

\begin{solution}
分 \(x<0\) 和 \(x\geqslant0\) 两种情况，有
\[
(1+x)(1-|x|)=
\begin{cases}
(1+x)^{2},&x<0,\\
(1+x)(1-x),&x\geqslant0.
\end{cases}
\]
当 \(x<0\) 时，结果除 \(x=-1\) 外均为正；当 \(x\geqslant0\) 时，因 \(1+x>0\)，还需 \(1-x>0\)，所以 \(0\leqslant x<1\). 合并得解集为 \(\{x\mid x<1\text{ 且 }x\neq-1\}\)，故选 D.
\end{solution}

\begin{problem}
函数 \( y = a^x \) 在 \([0,1]\) 上的最大值与最小值之和为 \(3\)，则 \( a = \)（\quad）

\choices
    {\(\frac{1}{2}\)}
    {\(2\)}
    {\(4\)}
    {\(\frac{1}{4}\)}
\end{problem}

\begin{answer}
B
\end{answer}

\begin{solution}
指数函数的底数满足 \(a>0\) 且 \(a\neq1\). 当 \(a>1\) 时，\(y=a^{x}\) 在 \([0,1]\) 上递增，最大值为 \(a\)，最小值为 \(1\)；当 \(0<a<1\) 时，最大值为 \(1\)，最小值为 \(a\). 两种情况下最大值与最小值之和都为 \(a+1\)，故 \(a+1=3\)，得 \(a=2\)，满足 \(a>1\)，故选 B.
\end{solution}

\begin{problem}
在 \((0,2\pi)\) 内，使 \(\sin x > \cos x\) 成立的 \(x\) 的取值范围是（\quad）

\choices
    {\(\left(\frac{\pi}{4},\frac{\pi}{2}\right)\cup \left(\pi ,\frac{5\pi}{4}\right)\)}
    {\(\left(\frac{\pi}{4},\pi\right)\)}
    {\(\left(\frac{\pi}{4},\frac{5\pi}{4}\right)\)}
    {\(\left(\frac{\pi}{4},\pi\right) \cup \left(\frac{5\pi}{4},\frac{3\pi}{2}\right)\)}
\end{problem}

\begin{answer}
C
\end{answer}

\begin{solution}
由 \(\sin x-\cos x=\sqrt{2}\sin\left(x-\frac{\pi}{4}\right)\)，原不等式等价于 \(\sin\left(x-\frac{\pi}{4}\right)>0\). 当 \(0<x<2\pi\) 时，只有
\[
0<x-\frac{\pi}{4}<\pi
\]
这一段满足正弦值为正，从而 \(\frac{\pi}{4}<x<\frac{5\pi}{4}\). 故选 C.
\end{solution}

\begin{problem}
设集合 \(M = \left\{x \mid x = \frac{k}{2} + \frac{1}{4}, k \in \Z\right\}\)，\(N = \left\{x \mid x = \frac{k}{4} + \frac{1}{2}, k \in \Z\right\}\)，则（\quad）

\choices
    {\(M = N\)}
    {\(M \subseteq N\)}
    {\(M \supseteq N\)}
    {\(M \cap N = \varnothing\)}
\end{problem}

\begin{answer}
B
\end{answer}

\begin{solution}
将集合中的数分别化为以 \(\frac14\) 为单位的形式：
\[
M=\left\{\frac{2k+1}{4}\mid k\in\Z\right\}.
\]
\[
N=\left\{\frac{k+2}{4}\mid k\in\Z\right\}=\left\{\frac{m}{4}\mid m\in\Z\right\}.
\]
因为 \(2k+1\) 是整数，所以 \(M\) 中每个元素都属于 \(N\)，即 \(M\subseteq N\)，故选 B.
\end{solution}

\begin{problem}
椭圆 \(5x^{2} + ky^{2} = 5\) 的一个焦点是 \((0,2)\)，那么 \(k =\)（\quad）

\choices
    {\(-1\)}
    {\(1\)}
    {\(\sqrt{5}\)}
    {\(-\sqrt{5}\)}
\end{problem}

\begin{answer}
B
\end{answer}

\begin{solution}
将方程除以 \(5\)，得 \(x^{2}+\frac{k}{5}y^{2}=1\). 因焦点在 \(y\) 轴上，椭圆的半长轴为 \(\sqrt{5/k}\)，半短轴为 \(1\)，焦半距满足
\[
c^{2}=\frac{5}{k}-1.
\]
已知焦点为 \((0,2)\)，故 \(c=2\)，于是 \(4=\frac{5}{k}-1\)，解得 \(k=1\)，故选 B.
\end{solution}

\begin{problem}
一个圆锥和一个半球有公共底面，如果圆锥的体积恰好与半球的体积相等，那么这个圆锥轴截面顶角的余弦值是（\quad）

\choices
    {\(\frac{3}{4}\)}
    {\(\frac{4}{5}\)}
    {\(\frac{3}{5}\)}
    {\(-\frac{3}{5}\)}
\end{problem}

\begin{answer}
C
\end{answer}

\begin{solution}
设圆锥和半球的公共底面半径为 \(r\)，圆锥高为 \(h\). 由体积相等，得
\[
\frac{1}{3}\pi r^{2}h=\frac{2}{3}\pi r^{3},
\]
所以 \(h=2r\). 设圆锥轴截面的顶角为 \(\alpha\)，将轴截面沿高分成两个直角三角形，则
\[
\tan\frac{\alpha}{2}=\frac{r}{h}=\frac12.
\]
因此
\[
\cos\alpha=\frac{1-\tan^{2}(\alpha/2)}{1+\tan^{2}(\alpha/2)}=\frac{1-1/4}{1+1/4}=\frac35.
\]
故选 C.
\end{solution}

\begin{problem}
已知 \(0 < x < y < a < 1\)，则有（\quad）

\choices
    {\(\log_a(xy) < 0\)}
    {\(0 < \log_a(xy) < 1\)}
    {\(1 < \log_a(xy) < 2\)}
    {\(\log_a(xy) > 2\)}
\end{problem}

\begin{answer}
D
\end{answer}

\begin{solution}
由 \(0<x<a\) 和 \(0<y<a\)，得 \(0<xy<a^{2}<1\). 因为 \(0<a<1\) 时函数 \(\log_{a}t\) 单调递减，所以
\[
\log_{a}(xy)>\log_{a}(a^{2})=2.
\]
故选 D.
\end{solution}

\begin{problem}
函数 \( y = x^{2} + bx + c \)（\(x \in [0, +\infty)\)）是单调函数的充要条件是（\quad）

\choices
    {\(b \geqslant 0\)}
    {\(b \leqslant 0\)}
    {\(b > 0\)}
    {\(b < 0\)}
\end{problem}

\begin{answer}
A
\end{answer}

\begin{solution}
函数的导数为 \(f'(x)=2x+b\). 当 \(b\geqslant0\) 时，对任意 \(x\geqslant0\) 都有 \(f'(x)\geqslant0\)，所以函数在 \([0,+\infty)\) 上单调递增. 当 \(b<0\) 时，\(f'(x)\) 在 \(x=-b/2\) 处变号，函数先减后增，不是单调函数. 因而充要条件是 \(b\geqslant0\)，故选 A.
\end{solution}

\begin{problem}
设 \(\theta \in \left(0, \frac{\pi}{4}\right)\)，则二次曲线 \(x^{2} \cot \theta - y^{2} \tan \theta = 1\) 的离心率的取值范围为（\quad）

\choices
    {\(\left(0, \frac{1}{2}\right)\)}
    {\(\left(\frac{1}{2}, \frac{\sqrt{2}}{2}\right)\)}
    {\(\left(\frac{\sqrt{2}}{2}, \sqrt{2}\right)\)}
    {\(\left(\sqrt{2}, +\infty\right)\)}
\end{problem}

\begin{answer}
D
\end{answer}

\begin{solution}
将方程写成标准形式
\[
\frac{x^{2}}{\tan\theta}-\frac{y^{2}}{\cot\theta}=1.
\]
故 \(a^{2}=\tan\theta\)，\(b^{2}=\cot\theta\)，双曲线的焦半距满足 \(c^{2}=a^{2}+b^{2}\). 于是离心率
\[
e=\frac{c}{a}=\sqrt{\frac{\tan\theta+\cot\theta}{\tan\theta}}=\sqrt{1+\cot^{2}\theta}=\csc\theta.
\]
由 \(0<\theta<\pi/4\)，有 \(0<\sin\theta<\sqrt{2}/2\)，所以 \(e>\sqrt{2}\). 当 \(\theta\to0^{+}\) 时 \(e\to+\infty\)，故取值范围为 \((\sqrt{2},+\infty)\)，选 D.
\end{solution}

\begin{problem}
从正方体的 \(6\) 个面中选取 \(3\) 个面，其中有 \(2\) 个面不相邻的选法共有（\quad）

\choices
    {\(8\) 种}
    {\(12\) 种}
    {\(16\) 种}
    {\(20\) 种}
\end{problem}

\begin{answer}
B
\end{answer}

\begin{solution}
正方体的 \(6\) 个面组成 \(3\) 对相对面. 要使所选的 \(3\) 个面中有一对不相邻的面，先从 \(3\) 对相对面中选出一对，再从剩下的 \(4\) 个面中选出第 \(3\) 个面. 因为只选 \(3\) 个面，不会同时选出两对相对面，所以没有重复计数，选法共有 \(3\times4=12\) 种，故选 B.
\end{solution}

\section{填空题}
\begin{problem}
据新华社 2002 年 3 月 12 日电，1985 年到 2000 年间，我国农村人均居住面积如图所示，
其中，从\fillinblank{}年\fillinblank{}年的五年间增长最快.

\bitmapfigure[max width=0.58\linewidth]{img/2002/syllabus_liberal/q13_fig1.png}
\end{problem}

\begin{answer}
\(1995\)，\(2000\)
\end{answer}

\begin{solution}
由图可知，三个五年段的人均居住面积增长量分别为 \(17.8-14.7=3.1\mathrm{~m^{2}}\)、\(21.0-17.8=3.2\mathrm{~m^{2}}\) 和 \(24.8-21.0=3.8\mathrm{~m^{2}}\). 最大的增长量出现在 \(1995\) 年至 \(2000\) 年，所以应填 \(1995\)，\(2000\).
\end{solution}

\begin{problem}
函数 \( y = \frac{2x}{1 + x} \)（\(x \in (-1, +\infty)\)）图象与其反函数图象的交点为\fillinblank{}.
\end{problem}

\begin{answer}
\((0,0)\)，\((1,1)\)
\end{answer}

\begin{solution}
记 \(f(x)=\frac{2x}{1+x}\). 在 \(x>-1\) 时，\(f'(x)=\frac{2}{(1+x)^{2}}>0\)，所以反函数存在. 函数图象与反函数图象的交点在直线 \(y=x\) 上，故令
\[
x=\frac{2x}{1+x}.
\]
由于 \(x>-1\)，可乘以 \(1+x\)，得 \(x(1+x)=2x\)，即 \(x(x-1)=0\)，所以 \(x=0\) 或 \(x=1\). 代回 \(y=x\)，交点为 \((0,0)\)，\((1,1)\).
\end{solution}

\begin{problem}
\((x^{2} + 1)(x - 2)^{7}\) 展开式中 \(x^{3}\) 的系数是\fillinblank{}.
\end{problem}

\begin{answer}
\(1008\)
\end{answer}

\begin{solution}
要得到展开式中的 \(x^{3}\) 项，第一部分的 \(x^{2}\) 必须乘以 \((x-2)^{7}\) 中的 \(x\) 项，第一部分的 \(1\) 必须乘以其中的 \(x^{3}\) 项. 由二项式定理，这两项的系数之和为
\[
\mathrm C_{7}^{1}(-2)^{6}+\mathrm C_{7}^{3}(-2)^{4}=7\times64+35\times16=448+560=1008.
\]
所以 \(x^{3}\) 的系数为 \(1008\).
\end{solution}

\begin{problem}
对于顶点在原点的抛物线，给出下列条件：

\begin{circlelist}
\item 焦点在 \(y\) 轴上；
\item 焦点在 \(x\) 轴上；
\item 抛物线上横坐标为 \(1\) 的点到焦点的距离等于 \(6\)；
\item 抛物线的通径的长为 \(5\)；
\item 由原点向过焦点的某条直线作垂线，垂足坐标为 \((2,1)\).
\end{circlelist}

能使这抛物线方程为 \(y^{2}=10x\) 的条件是\fillinblank{}. （要求填写合适条件的序号）
\end{problem}

\begin{answer}
\(\circled{2}\)、\(\circled{5}\)
\end{answer}

\begin{solution}
将目标抛物线写成 \(y^{2}=2px\) 的形式，得 \(p=5\)，所以焦点为 \(F\left(\frac52,0\right)\)，通径长为 \(2p=10\). 因此条件\(\circled{1}\)不成立，条件\(\circled{2}\)成立，条件\(\circled{4}\)不成立.

若横坐标为 \(1\)，则抛物线上的点满足 \(y^{2}=10\)，其到焦点的距离平方为
\[
\left(1-\frac52\right)^{2}+y^{2}=\frac94+10=\frac{49}{4},
\]
故距离为 \(\frac72\)，不是 \(6\)，条件\(\circled{3}\)不成立.

对于条件\(\circled{5}\)，点 \(H(2,1)\) 与焦点 \(F\left(\frac52,0\right)\) 确定的直线斜率为 \(-2\)，而 \(OH\) 的斜率为 \(\frac12\)，两斜率之积为 \(-1\)，所以 \(OH\perp FH\). 该直线过焦点，且原点到它的垂足为 \(H\)，条件\(\circled{5}\)成立. 因而合适条件的序号为 \(\circled{2}\)、\(\circled{5}\).
\end{solution}

\section{解答题}
\begin{problem}
如图，某地一天从 \(6\text{ 时}\) 至 \(14\text{ 时}\) 的温度变化曲线近似满足函数 \( y = A \sin(\omega x + \varphi) + b \).
\begin{enumerate}
\item 求这段时间的最大温差；
\item 写出这段时间的函数解析式.
\end{enumerate}

\bitmapfigure[max width=0.42\linewidth]{img/2002/syllabus_liberal/q17_fig1.png}
\FloatBarrier
\end{problem}

\begin{solution}
\begin{enumerate}
\item 由图可知，这段时间的最高温度为 \(30\mathrm{~^{\circ}C}\)，最低温度为 \(10\mathrm{~^{\circ}C}\)，所以最大温差为 \(30-10=20\mathrm{~^{\circ}C}\).
\item 由最高温度和最低温度得振幅 \(A=\frac{30-10}{2}=10\)，中线参数 \(b=\frac{30+10}{2}=20\). 从 \(x=6\) 时的最低点到 \(x=14\) 时的最高点经历半个周期，所以 \(\frac{T}{2}=14-6=8\)，从而 \(T=16\)，\(\omega=\frac{2\pi}{T}=\frac{\pi}{8}\). 取 \(x=6\) 时正弦值为 \(-1\)，有 \(\frac{\pi}{8}\times6+\varphi=\frac{3\pi}{2}\)，可取 \(\varphi=\frac{3\pi}{4}\). 因此这段时间的函数解析式为 \(y=10\sin\left(\frac{\pi}{8}x+\frac{3\pi}{4}\right)+20\)，其中 \(x\in[6,14]\).
\end{enumerate}
\end{solution}

\begin{problem}
甲、乙物体分别从相距 \(70\text{ 米}\) 的两处同时相向运动. 甲第 \(1\text{ 分钟}\) 走 \(2\text{ 米}\)，以后每分钟比前 \(1\text{ 分钟}\) 多走 \(1\text{ 米}\)，乙每分钟走 \(5\text{ 米}\).
\begin{enumerate}
\item 甲、乙开始运动后几分钟相遇？
\item 如果甲、乙到达对方起点后立即折返，甲继续每分钟比前 \(1\text{ 分钟}\) 多走 \(1\text{ 米}\)，乙继续每分钟走 \(5\text{ 米}\)，那么开始运动几分钟后第二相遇？
\end{enumerate}
\end{problem}

\begin{solution}
\begin{enumerate}
\item 甲运动 \(n\) 分钟所走的路程为 \(S_n=2+3+\cdots+(n+1)=\frac{n(n+3)}{2}\)，乙运动 \(n\) 分钟所走的路程为 \(5n\). 第一次相遇时两者路程之和为 \(70\)，所以
\[
\frac{n(n+3)}{2}+5n=70,
\]
即 \(n^{2}+13n-140=0\). 解得 \(n=7\) 或 \(n=-20\)，舍去负值，故第一次相遇在开始运动后 \(7\) 分钟.
\item 甲在第 \(10\) 分钟末走到 \(65\) 米，在第 \(11\) 分钟内到达乙的起点；乙在 \(14\) 分钟末到达甲的起点. 到 \(14\) 分钟时，甲按原方向累计走过
\[
S_{14}=\frac{14\times17}{2}=119\text{ 米}.
\]
其中 \(70\) 米用于到达乙的起点，折返后又走了 \(119-70=49\) 米，所以甲此时距甲的起点 \(70-49=21\) 米，而乙在甲的起点. 第 \(15\) 分钟内，甲的速度为 \(16\text{ 米/分}\)，乙的速度为 \(5\text{ 米/分}\)，两者相向而行，\(1\) 分钟后甲走到距甲起点 \(5\) 米处，乙也走到距甲起点 \(5\) 米处. 因而第二次相遇在开始运动后 \(15\) 分钟.
\end{enumerate}
\end{solution}

\begin{problem}
四棱锥 \(P - ABCD\) 的底面是边长为 \(a\) 的正方形，\(PB \perp\) 平面 \(ABCD\).
\begin{enumerate}
\item 若面 \(PAD\) 与面 \(ABCD\) 所成的二面角为 \(60^{\circ}\)，求这个四棱锥的体积；
\item 证明无论四棱锥的高怎样变化，面 \(PAD\) 与面 \(PCD\) 所成的二面角恒大于 \(90^{\circ}\).
\end{enumerate}

\bitmapfigure[max width=0.36\linewidth]{img/2002/syllabus_liberal/q19_fig1.png}
\FloatBarrier
\end{problem}

\begin{solution}
\begin{enumerate}
\item 建立空间直角坐标系，取 \(B=(0,0,0)\)，\(A=(a,0,0)\)，\(C=(0,a,0)\)，\(D=(a,a,0)\)，\(P=(0,0,h)\)，其中 \(h>0\). 平面 \(PAD\) 的一个法向量为 \(\bs{n}_1=(h,0,a)\)，底面的一个法向量为 \(\bs{n}_0=(0,0,1)\). 设二面角为 \(\alpha\)，则
\[
\cos\alpha=\frac{|\bs{n}_1\cdot\bs{n}_0|}{|\bs{n}_1||\bs{n}_0|}=\frac{a}{\sqrt{h^{2}+a^{2}}}=\cos60^{\circ}.
\]
所以 \(h=\sqrt3a\)，四棱锥体积为
\[
V=\frac13a^{2}h=\frac{\sqrt3}{3}a^{3}.
\]
\item 仍取上述坐标. 平面 \(PAD\) 和平面 \(PCD\) 可分别取外法向量 \(\bs{n}_1=(h,0,a)\)，\(\bs{n}_2=(0,h,a)\).
二者的长度均为 \(\sqrt{h^{2}+a^{2}}\)，且 \(\bs{n}_1\cdot\bs{n}_2=a^{2}>0\). 设两外法向量的夹角为 \(\beta\)，则 \(\cos\beta>0\)，从而 \(\beta<90^{\circ}\). 四棱锥内部所对应的二面角为 \(180^{\circ}-\beta\)，所以它恒大于 \(90^{\circ}\).
\end{enumerate}
\end{solution}

\begin{problem}
设函数 \(f(x) = x^2 + |x - 2| - 1\)（\(x \in \R\)）.
\begin{enumerate}
\item 判断 \(f(x)\) 的奇偶性；
\item 求函数 \(f(x)\) 的最小值.
\end{enumerate}
\end{problem}

\begin{solution}
\begin{enumerate}
\item \(f(0)=1\neq0\)，所以 \(f(x)\) 不是奇函数；又有 \(f(2)=3\)，\(f(-2)=7\)，即 \(f(-2)\neq f(2)\)，所以 \(f(x)\) 也不是偶函数.
\item 按 \(x<2\) 与 \(x\geqslant2\) 分段，得
\[
f(x)=
\begin{cases}
x^{2}-x+1=\left(x-\frac12\right)^{2}+\frac34,&x<2,\\
x^{2}+x-3,&x\geqslant2.
\end{cases}
\]
当 \(x<2\) 时，第一段在 \(x=\frac12\) 处取得最小值 \(\frac34\)；当 \(x\geqslant2\) 时，第二段递增，最小值为 \(f(2)=3\). 比较两段，函数的最小值为 \(\frac34\).
\end{enumerate}
\end{solution}

\begin{problem}
已知点 \(P\) 到两定点 \(M(-1,0)\)，\(N(1,0)\) 距离的比为 \(\sqrt{2}\)，点 \(N\) 到直线 \(PM\) 的距离为 \(1\)，求直线 \(PN\) 的方程.
\end{problem}

\begin{solution}
设 \(P=(x,y)\). 若直线 \(PM\) 垂直于 \(x\) 轴，则它是 \(x=-1\)，点 \(N\) 到它的距离为 \(2\)，与题设矛盾，因此可设直线 \(PM\) 的斜率为 \(k\). 由于 \(M=(-1,0)\)，直线 \(PM\) 的方程为 \(y=k(x+1)\). 点 \(N=(1,0)\) 到该直线的距离为 \(1\)，所以
\[
\frac{2|k|}{\sqrt{k^{2}+1}}=1,
\]
解得 \(k^{2}=\frac13\). 又由 \(PM=\sqrt2\,PN\)，有
\[
(x+1)^{2}+y^{2}=2\bigl((x-1)^{2}+y^{2}\bigr),
\]
即 \(x^{2}+y^{2}-6x+1=0\). 代入 \(y^{2}=\frac{(x+1)^{2}}{3}\)，得
\[
x^{2}+\frac{(x+1)^{2}}{3}-6x+1=0,
\]
化简为 \(x^{2}-4x+1=0\)，故 \(x=2+\sqrt3\) 或 \(x=2-\sqrt3\). 当 \(x=2+\sqrt3\) 时，\(|y|=1+\sqrt3=|x-1|\)；当 \(x=2-\sqrt3\) 时，\(|y|=\sqrt3-1=|x-1|\). 因而直线 \(PN\) 的斜率只能为 \(1\) 或 \(-1\)，且均经过 \(N=(1,0)\).
所以直线 \(PN\) 的方程为 \(y=x-1\) 或 \(y=-x+1\).
\end{solution}

\section{附加题}
\begin{problem}
\begin{enumerate}
\item 给出两块相同的正三角形纸片（如图 1，图 2），要求用其中一块剪拼成一个三棱锥模型，另一块剪拼成一个正三棱柱模型，使它们的全面积都与原三角形的面积相等，请设计一种剪拼方法，分别用虚线标示在图 1，图 2 中，并作简要说明；
\item 试比较你剪的正三棱锥与正三棱柱的体积的大小；
\item 如果给出的是一块任意三角形的纸片（如图 3），要求剪拼成一个直三棱柱，使它的全面积与给出的三角形的面积相等. 请设计一种剪拼方法，用虚线标示在图 3 中，并作简要说明.
\end{enumerate}

\bitmapfigure[max width=0.74\linewidth]{img/2002/syllabus_liberal/q22_fig1.png}
\end{problem}

\begin{solution}
\begin{enumerate}
\item 设原正三角形的边长为 \(2\)，面积为 \(S=\sqrt{3}\). 在图 1 中连接三边的中点，原三角形被分成四个边长为 \(1\) 的正三角形. 沿三条中位线折起三个角上的小三角形，使中间的小三角形作为底面，三个角上的小三角形作为侧面，即可拼成棱长为 \(1\) 的正三棱锥.

\item 设正三棱锥的底面边长和正三棱柱的底面边长均为 \(1\). 在图 2 的三个顶点处各剪出一个全等的四边形，使每个四边形沿原正三角形边的两条较长邻边长度均为原边长的 \(\frac14\)，并使其中一组对角为直角. 余下部分沿所设计的折线折起，成为缺少一个底面的正三棱柱；剪下的三个四边形拼成另一个边长为 \(1\) 的正三角形底面. 这样得到的正三棱柱两个底面都是边长为 \(1\) 的正三角形.

正三棱锥的底面面积为 \(\frac{\sqrt{3}}{4}\)，其高为
\[
h_{\text{锥}}=\sqrt{1-\left(\frac{\sqrt{3}}{3}\right)^{2}}=\frac{\sqrt{6}}{3}.
\]
正三棱柱的两个底面共占面积 \(\frac{\sqrt{3}}{2}\)，所以三个矩形侧面的总面积为 \(S-\frac{\sqrt{3}}{2}=\frac{\sqrt{3}}{2}\). 由于三个侧面的一边均为 \(1\)，设正三棱柱的高为 \(h_{\text{柱}}\)，由三个侧面的宽度均为 \(h_{\text{柱}}\)，得 \(3h_{\text{柱}}=\frac{\sqrt{3}}{2}\)，从而 \(h_{\text{柱}}=\frac{\sqrt{3}}{6}\).
因此 \(V_{\text{锥}}=\frac13\cdot\frac{\sqrt{3}}{4}\cdot\frac{\sqrt{6}}{3}=\frac{\sqrt{2}}{12}\). 又 \(V_{\text{柱}}=\frac{\sqrt{3}}{4}\cdot\frac{\sqrt{3}}{6}=\frac18\).
又因为 \(2\sqrt{2}<3\)，所以
\[
V_{\text{锥}}-V_{\text{柱}}=\frac{\sqrt{2}}{12}-\frac18=\frac{2\sqrt{2}-3}{24}<0,
\]
故 \(V_{\text{锥}}<V_{\text{柱}}\).

\item 设任意三角形为 \(\triangle ABC\)，内心为 \(I\). 分别取 \(AI\)，\(BI\)，\(CI\) 的中点 \(U\)，\(V\)，\(W\)，以 \(\triangle UVW\) 为直三棱柱的一个底面. 从 \(U\) 向 \(AB\)，\(AC\) 作垂线，从 \(V\) 向 \(BA\)，\(BC\) 作垂线，从 \(W\) 向 \(CA\)，\(CB\) 作垂线；沿这六条垂线剪开，得到顶点 \(A\)，\(B\)，\(C\) 附近的三个四边形. 将这三个四边形拼成与 \(\triangle UVW\) 全等的另一个底面，余下部分以 \(\triangle UVW\) 为一个底面，并沿相应折线折起成为直三棱柱的三个矩形侧面.

下面验证面积关系. 设 \(\triangle ABC\) 的面积为 \(S\)，内切圆半径为 \(r\). 设内切圆分别与 \(AB\)，\(BC\)，\(CA\) 相切，切点分别为 \(D\)，\(E\)，\(F\). 连接 \(D\)，\(E\)，\(F\) 与 \(I\)，则以 \(A\)，\(B\)，\(C\) 为顶点的三个四边形 \(ADIF\)，\(BDIE\)，\(CEIF\) 恰好分割 \(\triangle ABC\). 由于 \(U\) 是 \(AI\) 的中点，且 \(U\) 到 \(AB\)、\(AC\) 的垂足分别是 \(ADIF\) 在以 \(A\) 为中心、相似比为 \(\frac12\) 的相似变换下的像，同理可得另外两个顶点处的四边形也分别是相应四边形的相似比为 \(\frac12\) 的像. 因此，三个剪下的四边形面积之和为 \(\frac14S\). 又 \(U\)，\(V\)，\(W\) 分别是 \(AI\)，\(BI\)，\(CI\) 的中点，所以 \(\triangle UVW\) 是 \(\triangle ABC\) 关于 \(I\) 的相似比为 \(\frac12\) 的相似变换的像，从而
\[
S_{\triangle UVW}=\frac14S.
\]
三个剪下的四边形可以拼成 \(\triangle UVW\) 的另一个底面. 原三角形剩余部分的面积为 \(\frac34S\)，其中一个底面 \(\triangle UVW\) 占 \(\frac14S\)，所以三个矩形侧面的总面积为 \(\frac34S-\frac14S=\frac12S\). 因而折起后得到的直三棱柱全面积为
\[
2S_{\triangle UVW}+\frac12S=S,
\]
与原三角形面积相等.
\end{enumerate}
\end{solution}
